\section{Introduction}
\label{sec:introduction}

Learned routers are judged by the actions they select, but their attainable value
is bounded before learning begins. If every candidate fails on the same queries, a
perfect selector cannot improve utility over the best fixed action. If the only
successful action is already the least expensive, cost-aware selection has no
remaining trade-off to exploit. A router trained in either regime may be calibrated
and computationally efficient while solving an empty decision problem. We call the
missing prerequisite \emph{action-pool routability}: the pool must contain adequate
capability, query-dependent complementarity, measurable costs, and a feasible rule
for selecting among them.

This prerequisite is easy to overlook in retrieval-augmented generation (RAG).
RAG is not a single switch. A composite action can change the generator, sparse or
dense retrieval, second-hop expansion, external-memory organization, prompt length,
citation format, deterministic verification, and abstention policy. Recent systems
route among retrieval-augmented generators, retrievers, or direct/RAG/reasoning
pipelines \cite{ZhangEtAl2025RAGRouter,ZhaoEtAl2026R3AG,WangEtAl2026XRouter,
DingZhao2026RAGDiet}. These methods answer how to choose among useful options. They
do not remove the need to establish that the eligible options are useful and
complementary in the first place.

The issue is particularly acute for compact generators. Retrieved evidence helps
only when a model can attend to it, produce the requested schema, cite the exact
support, and abstain when support is absent. Model, task, context, and retrieval
characteristics jointly affect this behavior \cite{LiEtAl2025LaRA}; imperfect or
distracting context can reduce answer quality
\cite{OhThorne2023DetrimentalContexts,ParkLee2024RobustRALM}. Grounding also adds
work. Existing studies profile accuracy against tokens, latency, monetary cost, or
energy \cite{SchwartzEtAl2020GreenAI,HendersonEtAl2020EnergyReporting,
PanEtAl2025Electro,GuoEtAl2026RAGEnergy}. A scientific audit must therefore state
both its utility contract and its measurement boundary rather than treating
``quality'' or ``energy'' as interchangeable across systems.

We retrospectively apply an integrated failure-first audit to a controlled
diagnostic study of small-model RAG and memory actions. Adapting the established
single-best/virtual-best comparison from algorithm selection, the key diagnostic
compares the best fixed strict-success rate with a label-aware per-query oracle
over the same action pool. Their difference is routing headroom. A zero value
certifies that no selector can improve strict success on that matrix, regardless
of model architecture or calibration. This oracle is deliberately post hoc: it
diagnoses the opportunity present in sealed traces and is not reported as a
deployable policy \cite{Rice1976AlgorithmSelection,
LindauerEtAl2019AlgorithmSelection,LiEtAl2026LLMRouterBench}.

The empirical testbed contains 800 generated scenario groups and eight task classes
covering direct copying, document evidence, personal memory, composition, temporal
updates, authority conflicts, two-hop evidence, and deleted or missing facts. Five
pinned FP16 generators are compared through matched direct and composite-grounding
endpoints. Strict success requires answer correctness, exact support, valid
abstention, parseability, and exclusion of forbidden evidence. Physical telemetry
is restricted to selected-board GPU energy during generation; it is not a measure
of total RAG energy or carbon.

The study asks four questions:

\begin{enumerate}
    \item[\textbf{RQ1}] Does the resident-eligible action pool contain adequate capability
    and query-dependent routing headroom?
    \item[\textbf{RQ2}] When retrieved evidence is held constant, where does the
    pipeline lose utility between retrieval, generation, citation, and verification?
    \item[\textbf{RQ3}] How do generator and task class alter the success and
    GPU-generation-energy effect of the composite grounding intervention?
    \item[\textbf{RQ4}] Can the frozen learned controller satisfy its validation
    constraints, and does its sealed-test behavior agree with the pool audit?
\end{enumerate}

Across 11 resident-eligible actions, the best fixed route and the per-query oracle
both succeeded on 15 of 120 questions (12.5\%); every
success belonged to the same direct action. In contrast, the six offline reference
actions had 28.3 \pp{} of headroom, and the full 17-route matrix had 29.2 \pp{}.
Thus, complementarity existed in the experiment but was excluded from the learned
controller's action set. The remaining analyses explain that separation. All five
grounded endpoints received identical retrieved evidence, yet their evidence-task
success ranged from 0\% to 58.9\%. For the capable reference generators, choosing
direct operation on the two evidence-free task classes and grounding on the other
six improved strict success and used less measured generation energy than always
grounding. The learned controller, however, had no validation-feasible threshold
and reproduced a zero-success fallback on every test query.

The paper makes three contributions:

\begin{enumerate}
    \item an integrated failure-first protocol that adapts the established
    single-best/oracle gap to evidence-qualified composite RAG actions and combines
    it with interface, capability, physical-cost, and constraint-feasibility checks;
    \item a controlled retrieval-to-utilization analysis showing that identical
    evidence can yield sharply different schema, citation, answer, and strict-success
    outcomes across generators and task classes;
    \item a sealed negative case study that couples strict evidence-qualified
    utility with a precisely bounded physical GPU-energy measure and preserves
    validation infeasibility instead of repairing the policy after test access.
\end{enumerate}

The claim is not that a new routing algorithm outperforms prior work. It is that
routing opportunity is an empirical property of the eligible action pool and
should be audited before optimization. In this retrospective case, the audit
provides an impossibility certificate for strict-success improvement within the
resident-eligible matrix and identifies where utility complementarity remains.
